% !PRINTED 23
\weberpage{23}
Setzt man
\weq{9}{\frac{J}{\Delta^2} = 3z, \quad dz = \frac{1}{3}\frac{\Delta dJ-2Jd\Delta}{\Delta^3},}
so ergibt sich aus (8)
\[
K^2 = 27\Delta^6(4z^3+12Sz-T)
\]
und aus (7)
% [sic] The source prints 1/\sqrt{3}; equations (7)--(9) algebraically give \sqrt{3} (up to the branch sign).
\weq{10}{du = \frac{1}{\sqrt3}\frac{dz}{\sqrt{4z^3+12Sz-T}},}
und dies ist die \gesperrt{Weierstrasssche} Normalform für $g_2=-12S$,
$g_3=T$.

\begin{center}
\S\,7. \gesperrt{Elliptische Raumkurven vierter Ordnung.}
\end{center}

Man kann algebraische Funktionen einer Veränderlichen auch
durch Gleichungen zwischen mehreren Variablen darstellen, wenn
man die Anzahl der Gleichungen entsprechend vergrößert. Nimmt
man z.\,B. drei Variable $x$, $y$, $z$ und läßt zwischen ihnen zwei
Gleichungen
\weq{1}{\varphi(x,y,z) = 0, \qquad \psi(x,y,z) = 0}
bestehen, so kann man aus diesen beiden Gleichungen z.\,B. $x$
eliminieren und erhält eine Gleichung zwischen $y$ und $z$, durch
die $y$ als algebraische Funktion von $z$ definiert ist. Es kann
dann, wenn man gewisse Ausnahmefälle ausschließt, $x$ und jede
rationale Funktion von $x$, $y$, $z$ rational durch $y$ und $z$ dargestellt
werden. Die Integrale der Form
\weq{2}{\int F(x,y,z)\,dz,}
in denen $F$ eine rationale Funktion bedeutet, gehören dann zu
dem durch (1) dargestellten algebraischen Gebilde. Nimmt man
$x$, $y$, $z$ als Cartesische Koordinaten im Raume an, so stellt (1)
eine Raumkurve als den Durchschnitt zweier Flächen $\varphi=0$,
$\psi=0$ dar, und die Integrale (2) gehören zu dieser Raumkurve.
Die Kurve heißt wieder \gesperrt{elliptisch}, wenn diese Integrale elliptisch sind.

Wir wollen diese Betrachtungen auf die \gesperrt{Raumkurven vierter
Ordnung erster Spezies} anwenden, d.\,h. auf die Kurven vierter
Ordnung, die sich als vollständiger Durchschnitt zweier Flächen
zweiten Grades darstellen lassen.
